\section{The formal problem statement}\label{sec:problem-statement}

This section states the MONAD research problem.
\Cref{sec:introduction} onward develops MONAD-FORC~0.2 (Finite
Observer-Resource Causality), the first fixed-context profile of the
programme: the observer interface of \Cref{sec:observer} realizes a
resource-indexed probe family, and detectable intervention influence
in \Cref{sec:operational-causality} realizes a present-probe
distinction over interventions.

\subsection{Problem and scope}

MONAD is a pure mathematical research programme. MONAD does not start
from an AI system, a software architecture, a physical theory, or a
model of cognition. Applied systems, quantum theory, biology, and
information theory can provide examples and tests. They do not define
MONAD\@. The name MONAD is a project name. It does not denote a monad
in the sense of category theory~\cite{maclane1998categories}, although
the paper uses categories.

MONAD studies mathematical systems in which the set of meaningful
distinctions is not fixed. A system can:
\begin{enumerate}
  \item compress its present description;
  \item admit new distinctions that were not available in the old
    description;
  \item transform from one description to another;
  \item preserve some structure and lose other structure during that
    transformation;
  \item assign different relevance to the same structure along
    different trajectories;
  \item change the space in which later transformations occur.
\end{enumerate}

The theory must describe these operations with compatible, explicitly typed
constructions. A fixed ambient metaframework is permitted; its entire
vocabulary must not be confused with the system's currently admitted or
accessible distinctions. The theory must
not treat present indistinguishability as permanent identity. The
theory must not select a geometric, topological, probabilistic,
smooth, linear, or categorical substrate unless the problem forces
that structure. The first task is therefore not to find an equation.
The first task is to identify the correct mathematical objects and
operators.

\subsubsection{Core problem}

The shortest current statement is:

\Principle{MONAD seeks a calculus of endogenous distinguishability.}

``Distinguishability'' means that a mathematical description can
separate one object, behavior, relation, or continuation from
another. ``Endogenous'' means that the structure that determines those
distinctions can itself change inside the theory.

Compression works by deciding which distinctions can presently be
ignored. The difficulty begins when a later admissible question
depends on one of them:

\Principle{How can a finite description compress what is known without
silently destroying distinctions that an admissible future description
can require?}

The theory must also answer four further questions:
\begin{itemize}
  \item What does ``the same information'' mean when the language
    used to express that information changes?
  \item What does preservation mean when the space of possible
    distinctions itself changes?
  \item When does history matter even when two present descriptions
    appear equivalent?
  \item How can the theory detect that the present basis is
    insufficient, rather than only optimize inside that basis?
\end{itemize}

\subsubsection{What MONAD is not}

MONAD is not a theory of compression alone. It is not a theory of
dynamical systems in a fixed state space. It is not a theory of path
dependence alone. It is not ordinary category theory with different
names. It is not a claim about quantum mechanics.

MONAD does not assume retrocausality. It does not claim that a
physical future sends information into the present. It does not use
``chaos'' to mean that all possible outcomes physically coexist. It
does not assume that a tensor is the correct final object, that a
manifold is the correct substrate, or that every useful relation has a
scalar value. It does not assume that the four working postulates are
all primitive axioms. Some postulates can become theorems of deeper
structure. The final theory must discover which structure is primitive
and which structure is derived.

\subsection{The minimal mathematical starting point}

A category of description contexts records identities and associative
composition without requiring a metric, measure or linear structure.

Let $\Ctx$ be a category. An object $C \in \operatorname{Ob}(\Ctx)$ is
a description context. A morphism $u : C \to D$ is an admissible
transformation or extension between description contexts. This does
not mean that category theory is the final substrate. It only gives
identity and associative composition. These properties are already
required by the problem.

For each context $C$, let $\Carrier{C}$ be a carrier of entities
represented in that context. The carrier does not have to be a vector
space, a manifold, or finite. For each context $C$, let $\Probes{C}$
be the family of admissible probes. A probe is a map
$p : \Carrier{C} \to O_p$; the codomain $O_p$ can depend on the probe.
A probe represents one distinction that the current description can
express. For a morphism $u : C \to D$, assume, in the first model, a
corresponding transport $u_* : \Carrier{C} \to \Carrier{D}$. This
assumption can later be weakened.

% Lean module: MonadFormal.Sandbox
\subsubsection{Present and continuation-sensitive distinctions}
\label{ss:continuations}

For a context $C$, present indistinguishability is
\begin{equation}
  x \presentEq{C} y
  \iff p(x)=p(y)\quad\text{for every }p\in\Probes{C}.
  \label{eq:present-eq}
\end{equation}
This is equality of admitted observations, not equality of entities.
Let $\Conts{C}$ be an admitted family of continuations $k:C\to D$.
Their continuation-probe family is
\begin{equation}
  \ContProbes{C}=\{p\circ k_*:k\in\Conts{C},\;
    p\in\Probes{\operatorname{cod}(k)}\}.
  \label{eq:cont-probe-family}
\end{equation}
Define
\begin{equation}
  x\contEq{C}y\iff f(x)=f(y)
    \quad\text{for every }f\in\ContProbes{C}.
  \label{eq:cont-eq}
\end{equation}
The future enters through admitted tests of possible continuations;
no information from an actual future is assumed available.

Both relations use the same mathematical object: an indexed family
$F=(f_i:X\to O_i)_{i\in I}$, with possibly different codomains.
Write $x\approx_F y$ when $f_i(x)=f_i(y)$ for every $i$.
\leandef{MonadFormal.ProbeFamily,MonadFormal.indist}

\begin{theorem}[Observational equivalence]\label{thm:observational-equivalence}
For every probe family $F$, the relation $\approx_F$ is an equivalence
relation on $X$.
\end{theorem}
\begin{proof}
For each $i$, equality in $O_i$ is reflexive, symmetric and transitive.
Applying the corresponding equality law separately at every index proves
the three laws for $\approx_F$.
\leanproof{MonadFormal.indist_refl,MonadFormal.indist_symm,MonadFormal.indist_trans}
\end{proof}


\subsubsection{The change in viewpoint}

The continuation family, the probe family, the context, and the
carrier can all change, so the equivalence relation itself can change:
$\contEq{C} \to \contEq{C'}$. This is the formal meaning of a change
in the available space of distinctions. It is not only a new value
inside an old coordinate system. The coordinate system can become
insufficient.

\subsection{The four MONAD postulates}

The four postulates are working requirements, not yet final axioms.
The mature theory can show that one or more postulates follow from
deeper structure.

\subsubsection{Postulate I: generative openness}\label{ss:non-factorization}
A description can acquire a genuinely new distinction. An extension
$e : C \to C'$ is distinction-generating if there exist
$x, y \in \Carrier{C}$ with $x \contEq{C} y$ but
$e_*(x) \not\contEq{C'} e_*(y)$:

\Principle{An equivalence class in $C$ can split in $C'$.}

Such a split can merely reveal an old, previously unqueried bit; it is
not sufficient evidence of a new language or degree of freedom.
A second form of openness occurs when
the new carrier contains objects that are not transported old objects:
$\Carrier{C'} \setminus \operatorname{Im}(e_*) \neq \varnothing$.
Generative openness therefore has at least two forms, refinement of
old distinctions and extension beyond the old carrier. The theory
must keep them apart.

\paragraph{Generative openness as non-factorization.}

Let $q_C : \Carrier{C} \to Q_C$ be a compression. A new probe after
extension gives $f' = p' \circ e_*$. If $f' = \tilde f \circ q_C$ for
some $\tilde f$, the old compressed representation already contains
enough information for that new probe. If no such $\tilde f$ exists,
$f'$ does not factor through $q_C$. A new distinction is not merely a
new function. It is a function that cannot be expressed through the
old quotient. Define the extension defect set:
\begin{equation}
  \ExtDefect(e, q_C)
  = \{\, f' \in \ContProbes{C'} \circ e_* :
     f' \text{ does not factor through } q_C \,\}.\label{eq:defect}
\end{equation}
$\ExtDefect(e, q_C) = \varnothing$ means the old compression remains
sufficient under the extension. Otherwise the extension has exposed
distinctions that the old compression cannot support.

% Lean module: MonadFormal.Quotient
\subsubsection{Postulate II: anticipatory compression}\label{ss:anticipatory-compression}\label{ss:relative-compression}

The requirement is to preserve the distinctions admitted continuations
can demand. For a declared family, its exact set-based realization is a
sufficient representation.

A map $f:X\to O$ \emph{factors through} $q:X\to Y$ if there is
$\bar f:Y\to O$ with $f=\bar f\circ q$. A representation $q$ is
sufficient for $F$ if every probe factors through it:
\begin{equation}
  \forall i\in I\;\exists\bar f_i:Y\to O_i,
  \qquad f_i=\bar f_i\circ q.\label{eq:sufficient}
\end{equation}
\leandef{MonadFormal.Factors,MonadFormal.Sufficient}
This is exact preservation of declared distinctions. It is not, by itself,
minimal storage, efficient evaluation or a physically obtainable encoding.

\begin{theorem}[Coarsest sufficient quotient]\label{thm:coarsest-quotient}
Let $Q_F=X/{\approx_F}$ and $q_F(x)=[x]_F$. Then $q_F$ is
surjective and sufficient, and $q_F(x)=q_F(y)$ exactly when
$x\approx_F y$. For every sufficient surjection $q:X\to Y$, there
is a unique $h:Y\to Q_F$ such that $q_F=h\circ q$.
\end{theorem}
\begin{proof}
Define $\bar f_i([x]_F)=f_i(x)$. Observational equivalence makes
this independent of the representative, so each probe descends to
$Q_F$. Every quotient class has a representative. Conversely,
sufficiency implies that $q(x)=q(y)$ forces $f_i(x)=f_i(y)$ for
every $i$, hence $q_F(x)=q_F(y)$. Define $h(q(x))=q_F(x)$.
The preceding implication makes this well-defined, and surjectivity
of $q$ makes it both everywhere defined and unique.
\leanproof{MonadFormal.toQuot_sufficient,MonadFormal.toQuot_eq_iff,MonadFormal.toQuot_surjective,MonadFormal.quotient_universal}
\end{proof}

\begin{theorem}[Unique descent and relabeling]\label{thm:unique-descent}
If $q:X\to Y$ is surjective and $f:X\to O$ is constant on its
fibers, then $f=\bar f\circ q$ for a unique $\bar f:Y\to O$.
If surjections $q:X\to Y$ and $r:X\to Z$ have the same fibers,
their values differ only by a bijection $h:Y\to Z$ satisfying
$r=h\circ q$.
\end{theorem}
\begin{proof}
For each $y$, choose $x$ with $q(x)=y$ and set $\bar f(y)=f(x)$.
Fiber constancy removes dependence on the choice; evaluating any
other candidate at $q(x)$ proves uniqueness. Apply this construction
to $r$ through $q$, and to $q$ through $r$. The two resulting maps
compose to the identity on the images, which are the whole codomains.
\leanproof{MonadFormal.unique_descent,MonadFormal.same_fibers_relabeling}
\end{proof}

No finiteness hypothesis is needed. The construction can be
noncomputable: existence of a quotient is not an algorithm for
recognizing its classes.


\paragraph{Compression is relative, not absolute.}

There is no MONAD notion of lossless compression without a declared
distinction family. Lossless means lossless relative to
$\ContProbes{C}$. If $\Conts{1} \subseteq \Conts{2}$ then
$x \contEq{\Conts{2}} y \implies x \contEq{\Conts{1}} y$. A larger
continuation family can only preserve or refine the equivalence
relation.

\Principle{An expansion of relevant continuations can force a finer
compression.}

\paragraph{The universal compression obstruction.}

% Lean module: MonadFormal.NoGo
\begin{theorem}[Universal compression obstruction]\label[theorem]{prop:no-go}
Let $F$ jointly separate $X$: for any $x\ne y$, some probe has
$f_i(x)\ne f_i(y)$. Every representation sufficient for $F$ is
injective. In particular, the conclusion applies when the probes
are composites of admitted future extensions and subsequent tests.
\end{theorem}
\begin{proof}
Suppose $q(x)=q(y)$. For every $i$, sufficiency supplies
$f_i=\bar f_i\circ q$, so $f_i(x)=f_i(y)$. Joint separation
therefore forces $x=y$.
\leanproof{MonadFormal.no_go}
\end{proof}


\Principle{Unrestricted future openness plus perfect future-safe
compression implies no nontrivial compression.}

The proposition leaves one perfectly safe strategy: retain enough to
answer every admissible future probe. Under joint separation, that
strategy retains everything. Perfect caution has solved the problem
by refusing to compress. Nontrivial compression therefore requires
giving something up: which continuations count, whether the
representation may be revised, how much approximation is tolerated,
how much uncertainty or explicit defect is carried, what can be
recovered afterward, what resources bound the observer, and how the
accounts compose. The later sections study these relaxations
separately. This tension
is structural. It is not an implementation problem.

\paragraph{Approximate anticipatory compression.}

Exact sufficiency can be too strong. A later theory can add a measure
$\mu_C$ over admissible continuations and a distortion
$\ell(q, \gamma)$, giving a resource-sensitive compression:
\begin{equation}
  q^{\star} = \arg\min_q \left[ \operatorname{Complexity}(q)
  + \beta\, \mathbb{E}_{\gamma \sim \mu_C}\, \ell(q, \gamma) \right].\label{eq:approx-compression}
\end{equation}
This resembles rate-distortion theory~\cite{shannon1959coding} and
information bottleneck methods~\cite{tishby1999information}, and the
complexity term connects to minimum description length and Kolmogorov
complexity~\cite{rissanen1978modeling,li2008introduction}. However,
the continuation family itself can change, so
$\mu_C$ can be incomplete. Low expected loss under the current
continuation model is not structural sufficiency under unknown
extensions, and the two must be kept apart.

\subsubsection{Postulate III: invariant accounting}

A transformation must not only produce a result. It must admit an
account of what happened to distinguishability. A transformation
$T : C \to D$ can do at least four qualitatively different things.
\emph{Preservation}: a distinction in $C$ remains a distinction in
$D$. \emph{Hiding}: a distinction is not visible under the current
target probes, but an admissible recovery can reconstruct it.
\emph{Collapse}: a distinction in $C$ cannot be recovered after
transformation. \emph{Novelty}: the target contains distinctions that
cannot be generated from the source representation. The four cases
must not merge into one generic notion of information loss; the
accounting exists to keep them apart.

\paragraph{Distinction collapse.}

A first definition:
\begin{equation}
  \operatorname{Coll}(T)
  = \{\, (x, y) : x \not\contEq{C} y,\;
     T_*(x) \contEq{D} T_*(y) \,\}.\label{eq:collapse}
\end{equation}
This does not yet determine whether the distinction is irreversibly
destroyed. For that, MONAD needs recovery.

\paragraph{Recoverability.}

Assume $T$ induces $\bar T : \Carrier{C}/{\contEq{C}} \to
\Carrier{D}/{\contEq{D}}$. A distinction is recoverable on a specified
substructure if there is a recovery map $R$ with
$R \circ \bar T = \operatorname{id}$ on that substructure. Loss of
immediate visibility does not imply destruction. MONAD must separate
\emph{inaccessible} from \emph{irrecoverable}. Quantum information
theory gives a useful existing
model~\cite{petz1986sufficient,hayden2004structure}.

\paragraph{Novelty under transformation.}

In a general category one cannot write $D - \operatorname{Im}(T)$
without additional structure, so novelty must first be defined
intrinsically: a target distinction is novel relative to $T$ if it
does not arise from transported source distinctions under the
admissible transport rules. If an exact or additive structure is later
justified, the concept sharpens through image, kernel, cokernel, exact
sequences, extension classes, and cohomological obstructions. MONAD
must earn these tools, not assume them for convenience.

\paragraph{Invariant transport.}

For a structure $I_C$ regarded as invariant in $C$, MONAD needs a
transport operation of the form
\begin{equation}
  \TransportOp_T(I_C) = (I_D,\; L_T,\; N_T,\; \Pi_T),\label{eq:transport}
\end{equation}
where $I_D$ is the transported structure, $L_T$ records explicit loss,
$N_T$ records newly expressible structure, and $\Pi_T$ records the
witness for the transport claim. The domain and codomain of
$\TransportOp_T$ are deliberately untyped at this stage. Whether $I_C$
is a predicate, a subobject, an equivalence relation, a conserved
quantity, or a theory is an open typing problem of the programme, and
it is prior to any composition law: the class of invariant objects
must be fixed before an interaction term can be posed. The notation is
a first schema. The requirement is durable:

\Principle{Loss and novelty must be first-class outputs of
transformation. They must not be hidden in prose.}

\paragraph{Composition creates the accounting problem.}

For $C \xrightarrow{f} D \xrightarrow{g} E$, one can account for
$g \circ f$ directly, or account for $f$ and then $g$. These accounts
need not agree trivially. MONAD must determine whether
$\Acct(g \circ f)$ is completely determined by $\Acct(f)$ and
$\Acct(g)$, or whether an irreducible interaction term exists. The
programme states this as an open question, not as an equation:

\Principle{Open problem. Under which substrates does transformation
accounting admit functorial composition, and what obstruction measures
the failure of such composition?}

If composition requires an independent obstruction object, the theory
must identify its natural home (a cocycle, an extension class,
curvature, a derived functor, or nothing at all) and must not select
one before the substrate requires it. An equation with an interaction
term would commit the theory to an additive picture it has not yet
earned.

\subsubsection{Postulate IV: trajectory-sensitive relevance}\label{ss:trajectory-relevance}

A distinction can have different importance along different
trajectories. This is stronger than ordinary path dependence: the
relevance itself can depend on the path. For a path
$\gamma : C_0 \to C_1 \to \cdots \to C_n$ and a family $\Gamma$ of
admissible continuations from $C_n$, MONAD requires a relevance object
$\Relevance(\gamma, \Gamma)$. The letter $\Gamma$ names a family
taken from a path endpoint; $\Conts{C}$ names the family attached to
one context. The theory must not assume that
$\Relevance$ is a scalar. It can be a relation, an order, a subobject,
a family of probes, an operator, a multilinear object, or another
structure forced by the final substrate. A scalar attention weight is
only one special case.

\Principle{Open problem. Determine the codomain and type of the
trajectory-relevance object $\Relevance$.}

Until that type is fixed, $\Relevance$ is a requirements placeholder,
not a mathematical object. Postulates I and II are at present more
formal than III and IV, so the four pillars are not yet mathematically
symmetric. The primitive-versus-derived question below already asks
whether relevance is primitive; that reduction, from relevance to
continuation structure, should be investigated before more machinery
is added.

\paragraph{Relevance depends on trajectory, not only endpoint.}

For two paths $\gamma_1, \gamma_2 : C_0 \leadsto C_n$ with the same
endpoints, MONAD permits
$\Relevance(\gamma_1, \Gamma) \neq \Relevance(\gamma_2, \Gamma)$:

\Principle{The same endpoint can have different applicability
structure because the histories differ.}

This is the exact mathematical issue behind the holonomy analogy.
MONAD does not assume a connection or a bundle. If later structure
permits a transport operator $\TransportOp_\gamma$, a closed loop can
have $\TransportOp_\gamma \neq \operatorname{id}$, a holonomy-like
defect. The theory must first earn the structure that makes the
statement valid.

\paragraph{Past and continuation jointly condition the present.}

The current distinguishability relation can depend on both the path by
which the context was reached and the admissible continuations from
it, giving the more complete relation
$x \contEq{C;\,\gamma,\,\Gamma} y$. The path represents historical
dependence; the continuation family represents anticipatory
dependence. No backward physical causation is required. The present
mathematical structure is indexed by both sides of a path.

\paragraph{Basis discovery.}

Ordinary relevance systems ask which represented dimensions are
important. MONAD must also permit the stronger question: is the
current family of dimensions itself insufficient? The formal witness
is again non-factorization (\Cref{ss:non-factorization}): if a
trajectory makes a probe $p'$ relevant and $p'$ does not factor
through the current quotient, the representation cannot express a
distinction the trajectory requires.
Trajectory-sensitive relevance therefore triggers generative openness.
The four postulates cannot remain independent modules.

\subsection{Compatibility between the four postulates}

Compatibility requires exact translations and shared assumptions between
the four pillars, whether their mathematics is reused or extended.
The subsections below state the minimum cross-postulate laws. The laws
in this group are definitional sanity checks: each follows quickly
from the definitions it relates. They validate the vocabulary. They
are useful whether or not a new calculus is needed.

% Lean module: MonadFormal.Defect
\subsubsection{Compression-openness compatibility}

\begin{theorem}[Empty extension defect]\label{thm:empty-defect}
For a target probe family $F'$ and transport $e:X\to X'$, the
extension defect $\ExtDefect(e,q)$ is empty if and only if $q$ is
sufficient for the pulled-back family $(f'_i\circ e)_i$.
\end{theorem}
\begin{proof}
By definition a probe belongs to the defect exactly when its
pullback fails to factor through $q$. Emptiness says that every
pullback factors through $q$, which is the definition of sufficiency.
\leanproof{MonadFormal.defectEmpty_iff_sufficient}
\end{proof}
This equivalence is a diagnostic baseline, not a construction of an
obtainable refinement or a bound on its cost.

% Lean module: MonadFormal.Monotonicity
\subsubsection{Continuation monotonicity}\label{ss:monotonicity}

\begin{theorem}[Refinement under family growth]\label{thm:family-growth}
Let $F$ and $G$ be probe families on the same carrier. If $x$ and
$y$ agree on every probe in $F\cup G$, they agree on every probe
in $F$. Thus enlarging a continuation-probe family can only refine
its observational equivalence.
\end{theorem}
\begin{proof}
Each probe of $F$ is a member of the enlarged family. Restrict its
universal equality condition to those members. The formal family uses
a disjoint union of indices; duplicate probes do not change the relation.
\leanproof{MonadFormal.indist_mono}
\end{proof}

% Lean module: MonadFormal.Extension
\subsubsection{Covered extensions and irreversible loss}

\begin{theorem}[Fiber-constancy criterion]\label{thm:fiber-constancy}
For a surjection $q:X\to Y$, a probe $f:X\to O$ factors through
$q$ if and only if $q(x)=q(y)$ implies $f(x)=f(y)$.
\end{theorem}
\begin{proof}
If $f=\bar f\circ q$, apply $\bar f$ to $q(x)=q(y)$.
Conversely, choose a representative $x_y$ in each fiber and define
$\bar f(y)=f(x_y)$. Constancy on fibers gives
$\bar f(q(x))=f(x)$. Surjectivity avoids an unspoken assumption
that values can be assigned on unused codomain elements.
\leanproof{MonadFormal.factors_iff_fiber_const}
\end{proof}

\begin{theorem}[Closure of covered extensions]\label{thm:covered-extension}
Suppose every probe $f'_i\circ e$ pulled back along $e:X\to X'$
factors through the old quotient $q_F$. Then
$x\approx_F y$ implies $e(x)\approx_{F'}e(y)$.
\end{theorem}
\begin{proof}
For each $i$, write $f'_i\circ e=g_i\circ q_F$.
Observational equivalence gives $q_F(x)=q_F(y)$; applying $g_i$
gives $f'_i(e(x))=f'_i(e(y))$.
\leanproof{MonadFormal.indist_pullback_of_factors}
\end{proof}
Selecting another member of an already covered family cannot split
an old class. Family enlargement and selection within a fixed family
are different context changes.

\begin{theorem}[Non-recovery from an identified pair]\label{thm:no-recovery}
If $q(x)=q(y)$ but $q'(e(x))\ne q'(e(y))$, there is no
$r:Y\to Y'$ with $r\circ q=q'\circ e$.
\end{theorem}
\begin{proof}
Such an $r$ would give
$q'(e(x))=r(q(x))=r(q(y))=q'(e(y))$, contradicting the hypothesis.
\leanproof{MonadFormal.no_recovery}
\end{proof}
This is a statement about recovery from $q$ alone. It does not exclude
an archive, a fresh measurement, or side information that distinguishes
the two states.

% Lean module: MonadFormal.Refinement
\subsubsection{Semantic refinement versus obtainable refinement}

Let $q:X\to Y$ be retained information and $r:X\to Z$ a newly
required distinction. Define $q^+(x)=(q(x),r(x))$.
\leandef{MonadFormal.refine}

\begin{theorem}[Least common refinement]\label{thm:least-refinement}
The map $q^+$ preserves both $q$ and $r$, and
\[
 q^+(x)=q^+(y)\iff q(x)=q(y)\ \text{and}\ r(x)=r(y).
\]
If both $q$ and $r$ factor through $s:X\to W$, then so does $q^+$.
Thus $q^+$ is least in the order determined by factorization.
\end{theorem}
\begin{proof}
The coordinate projections recover $q$ and $r$. Equality of ordered
pairs is equality of both coordinates. If $q=g\circ s$ and
$r=h\circ s$, then $q^+=(g,h)\circ s$.
\leanproof{MonadFormal.refine_fibers,MonadFormal.refine_preserves,MonadFormal.refine_least}
\end{proof}

\begin{theorem}[Obtainability from declared access]\label{thm:obtainable-refinement}
Let $A:X\to W$ describe the information actually accessible to the
system. The refined value $q^+$ factors through $A$ if and only if
both $q$ and $r$ factor through $A$. Moreover, deterministic
postprocessing of $A$ cannot separate any pair that $A$ identifies.
\end{theorem}
\begin{proof}
The forward implication composes the two projections with a decoder
for $q^+$. The reverse implication is the least-refinement property.
For the last statement, $A(x)=A(y)$ implies $h(A(x))=h(A(y))$ for
every function $h$.
\leanproof{MonadFormal.factors_trans,MonadFormal.refine_obtainable_iff,MonadFormal.processing_preserves_fibers}
\end{proof}
The definition of $q^+$ supplies no access to $r(x)$. In particular,
the analyst's access to $x$ is not an admissible operation of the system.

\begin{example}[New carrier states need not split old states]
\label{ex:carrier-extension}
Let $X=\{*\}$, $X'=\{*\}\times\{0,1\}$ and $e(*)=(*,0)$.
The second-coordinate probe is constant after pullback to $X$, so
it factors through the constant old representation. Nevertheless,
it separates $(*,0)$ and $(*,1)$ in $X'$, and $(*,1)$ is outside
the image of $e$.
\end{example}
\begin{proof}
The pulled-back probe has value zero. The two displayed target values
are zero and one; no image point has second coordinate one.
\leanproof{MonadFormal.extra_bit_pullback_factors,MonadFormal.extra_bit_separates_new_states,MonadFormal.extra_state_outside_image}
\end{proof}
Pullback sufficiency over old states therefore does not describe all
distinctions of an enlarged carrier.

% Lean module: MonadFormal.Counterexample
\subsubsection{Revelation without a new carrier}

\begin{example}[An old bit becomes relevant]\label{ex:old-bit}
Take $X=\{0,1\}$ and initially admit no probes. The constant map
$q:X\to\{*\}$ is sufficient. Add the identity probe on the same
carrier. The old class splits, and $q$ is no longer sufficient.
Every representation sufficient for the identity probe is injective.
\end{example}
\begin{proof}
The empty family imposes no equalities. After extension, the identity
probe returns zero on one state and one on the other. Any decoder of
the constant value would have to return both, which is impossible.
The identity probe jointly separates the carrier, so the final claim
is the universal compression obstruction.
\leanproof{MonadFormal.base_indist,MonadFormal.extended_split,MonadFormal.const_sufficient_before,MonadFormal.const_insufficient_after,MonadFormal.idFamily_separating,MonadFormal.no_go_on_bool}
\end{proof}
Non-factorization establishes insufficiency relative to the changed
family. It does not establish that a new physical degree of freedom,
language or carrier element came into existence.


\subsubsection{Compression-invariant compatibility}
An invariant claim is prohibited when the invariant does not factor
through the compressed representation: no invariant claim may depend
on a distinction already removed by compression. This connects
Postulates II and III\@.

\subsubsection{Path-accounting compatibility}
Transport along concatenated paths must compose,
$\TransportOp_{\gamma_2 * \gamma_1} = \TransportOp_{\gamma_2} \circ
\TransportOp_{\gamma_1}$, unless the theory explicitly records a
defect that explains the failure. If closed-loop transport is
nontrivial, endpoint-only accounting is insufficient. This connects
Postulates III and IV\@.

\subsubsection{Relevance-openness compatibility}
If relevance identifies a distinction not expressible in the current
quotient, it cannot be silently discarded: the representation
refines, the context extends, the distinction is declared
unresolved, or a controlled approximation is recorded. Relevance
cannot create a requirement and then allow the representation to
ignore it silently. This connects Postulates IV and I.

\subsubsection{Open-world honesty}

Current indistinguishability is not equality. Current non-relevance is
not permanent irrelevance. Failure to produce a distinction is not
proof that the distinction does not exist. MONAD therefore requires an
explicit unresolved status: when $x \presentEq{C} y$ but an admissible
continuation separates them, any claim requiring their permanent
identity is unsupported, and the correct status is
\emph{underdetermined}. Confusion here does not mean random behavior.
It means the present distinction structure is insufficient to justify
a stronger claim.

\subsubsection{The wrong-basis problem}

A system can be optimal inside a bad representation. In a smooth
specialization, $\nabla_M L(x^{\star}) = 0$ can hold while an
extension $i : M \hookrightarrow M'$ admits a direction
$v \in T_{i(x^{\star})} M'$ with $D_v L < 0$; the optimum was relative
to the old space. The general condition needs no derivatives:

\Principle{Minimality in $C$ does not imply minimality after an
admissible extension $C \to C'$.}

MONAD must distinguish optimization failure from representation
failure.

\subsection{Choosing the mathematical substrate}

Each mathematical structure gives useful operators and imports
assumptions. Differentiation needs local limit structure. A Riemannian
metric needs a smooth manifold with inner products. Kernels and
cokernels need categorical structure. Exact sequences need more.
Relative entropy needs probabilistic or operator-algebraic structure.
Holonomy normally needs a connection. Cohomology needs a suitable
complex, site, or sheaf.

\Principle{Add mathematical structure only when a MONAD requirement
forces it.}

This is the substrate principle.

\subsubsection{Minimal first sandbox}\label{ss:sandbox}

The first formal sandbox is deliberately weak. Use sets as carriers,
functions as transformations, families of probes and continuations,
equivalence relations, and quotient sets. The sandbox suffices to test
generative openness, continuation-sensitive equivalence, exact
compression, factorization failure, the no-go result, and simple
composition questions.

If the model becomes insufficient, the failure itself identifies what
structure to add:
\begin{itemize}
  \item an order, metric, divergence, or measure for quantitative loss;
  \item bundle or categorical structure for path transport;
  \item exact or abelian categories for algebraic accounting;
  \item probability for anticipation;
  \item smooth structure for sensitivity.
\end{itemize}

\subsubsection{Primitive versus derived structure}

The theory must search for the smallest axiom set
$\mathcal{A} = \{A_1, \ldots, A_n\}$ with
$\mathcal{A} \vdash P_1, P_2, P_3, P_4$ and the required coherence
laws, minimizing $|\mathcal{A}|$. A postulate that follows from the
other structure must become a theorem, not remain an independent
axiom.

For every proposed object, ask whether it is primitive or can be
derived from a smaller structure. Does relevance follow from
continuation structure? Does openness follow from failure of
factorization? Does invariant accounting arise from a universal
property? Does path dependence arise from non-flat transport? Can
dimension be defined through separating families of probes? The theory
must prefer derivation over declaration.

\subsubsection{The intrinsic-geometry requirement}

A mature theory must be representation-invariant: for an isomorphic
re-description $\phi : C \to C'$, a valid MONAD quantity transforms
naturally, independent of names and presentation. The long-term goal
is intrinsic operators that make transformation properties directly
visible, in the role the determinant plays for oriented volume.

\subsubsection{A provisional transformation signature}\label{ss:signature}

One provisional form is
\begin{equation}
  \Sigma(T) = (P_T,\; C_T,\; R_T,\; N_T,\; H_T),\label{eq:signature}
\end{equation}
recording preserved structure, collapsed distinctions,
recoverability, novelty, and path dependence. The tuple is a
requirement on what the final mathematics must expose. One deeper
object from which these quantities follow would be preferable.

\subsubsection{One structure, not four}

The goal is one structure in which the four postulates are natural
aspects of the same object. Openness changes the quotient structure.
Compression changes available distinctions. Invariant transport
constrains admissible transformations. Path structure changes
relevance. Relevance can force new distinctions. New distinctions
invalidate old compressions. This is a research criterion, not yet a
result.

\subsection{Existing mathematics that MONAD must reuse}

MONAD must not reinvent results that already exist. Existing
mathematics must be imported where the match is exact. A gap must be
stated precisely where the match is incomplete.

\subsubsection{Partition and quotient theory}
This gives indistinguishability classes, coarse-graining, refinement,
and quotient representations. A fixed carrier produces a lattice of
partitions. Changing these data is not by itself a gap: compare maps
between carriers, partition refinements and the exact factorization
results before introducing another construction.

\subsubsection{Coalgebra and behavioral equivalence}\label{sss:coalgebra}
Bisimulation and observational
equivalence~\cite{rutten2000universal,jacobs2016introduction} come
close to $x \contEq{C} y$. The comparison must specify the observation
signature and transition model. Observational equivalence and bisimulation
are not interchangeable without their respective hypotheses; changes of
signature must be compared with the corresponding established machinery.

\subsubsection{Abstract interpretation}\label{sss:abstract-interpretation}
Sound approximation and refinement~\cite{cousot1977abstract} are
directly relevant to safe compression, and completeness of an
abstract domain makes the insufficiency of an abstraction itself a
formal object~\cite{giacobazzi2000making}. A changing domain is not
an escape from abstraction refinement. The exact partition correspondence
is proved in \Cref{thm:partition-completeness,cor:budget-completeness}.
The separate operational question is whether the refined information can
be obtained from actual retained access, and at what cost.

\subsubsection{Sufficient statistics and information theory}
Sufficiency, rate-distortion, and information
bottleneck~\cite{cover2006elements,shannon1959coding,tishby1999information}
approximate anticipatory compression. Computational mechanics
sharpens the pattern: causal states are the minimal sufficient
statistic of history for predicting
futures~\cite{shalizi2001computational}. For a later change of relevance,
compare fixed-target, multitask and uncertain-target formulations. A
representation sufficient for one target need not be sufficient for a
different one; naming that change does not itself establish a new theorem.

\subsubsection{Quantum information theory}\label{ss:qit}
% monad-open: quantum-correspondence

For a quantum channel $\Phi$ the data processing inequality
holds~\cite{nielsen2010quantum}:
\begin{equation}
  D(\rho \Vert \sigma) \ge D(\Phi(\rho) \Vert \Phi(\sigma)).\label{eq:dpi}
\end{equation}
Equality conditions connect to Petz-type recovery: sufficiency of a
channel is characterized by the existence of a recovery
map~\cite{petz1986sufficient,hayden2004structure}. This is directly
relevant to the difference between hidden and irrecoverable. The precise
state, channel, support and recovery hypotheses must travel with the
imported theorem. The MONAD comparison question is:

\Principle{What happens to sufficiency and data-processing-like
bounds when the family of admissible distinctions itself expands?}

\subsubsection{Renormalization and coarse-graining}

Renormalization~\cite{wilson1975renormalization} disciplines
scale-dependent description, effective
degrees of freedom, and relevant operators. Determine which changing
descriptions already fit a renormalization construction. Neither a
shared coarse-graining vocabulary nor a fixed microscopic theory proves
equivalence or inequivalence with MONAD's operational problem.

\subsubsection{Contextuality and sheaf methods}

These handle context-dependent observability, local compatibility, and
local-to-global obstructions~\cite{abramsky2011sheaf}. An extensible
context family needs an explicit comparison of covers and their maps;
the fact that an individual theorem fixes a cover is not a demonstrated
limitation of the framework.

\subsubsection{Differential geometry, connections, and holonomy}\label{ss:holonomy}

These supply transport, path dependence, curvature, and
holonomy~\cite{kobayashi1963foundations}: a
close model for \Cref{ss:trajectory-relevance}. The manifold, bundle
and connection are additional hypotheses. Holonomy is a candidate specialization of
trajectory sensitivity, not a MONAD axiom.

\subsubsection{Exact and homological algebra}\label{ss:homological}

Kernel, image, cokernel, exact sequences, and cohomology offer a
canonical accounting language~\cite{weibel1994introduction}. Use it
when additive or exact structure is actually available; do not introduce
such structure solely to make an accounting metaphor look formal.

\subsubsection{Fibrations, indexed categories, and institutions}

These permit different contexts to carry different local languages,
with transport between
fibers~\cite{goguen1992institutions,jacobs1999categorical}. They are
strong candidates for the
architecture of variable description spaces. Existing frameworks already
describe language families and their transformations; this statement also
uses an ambient metaframework. The task is to translate the operational
problem, compose applicable results and record exact unmet hypotheses.
No absence-of-a-single-framework claim supplies evidence of a gap.

\subsubsection{Variational and path-space mathematics}

A functional over complete trajectories models whole-path constraint
without backward causation. Changes of admissible path families must be
tested against ordinary constrained optimization, control and state
augmentation, with their memory and computational costs stated.

\subsubsection{Open-ended evolution and adaptive systems}

These fields study new structures, changing networks, and changing
state spaces at the application level~\cite{bedau2000open}. A concrete
update rule can be useful mathematics, not a defect to reject for being
domain-specific. Its assumptions must be separated from any proposed
general law of changing distinguishability.

\subsubsection{Information geometry}

Divergences, metric structure, and dual connections give intrinsic
geometry to statistical distinguishability~\cite{amari2000methods}. This becomes useful if a
probabilistic substrate is earned. A changed statistical model needs a
specified comparison map and regularity conditions; it is not automatically
a new geometry or an obstruction to known information geometry.

\paragraph{Limits of the quantum analogy.}

MONAD must not equate wave-function collapse with ontology change.
The imported lessons are narrower:
\begin{itemize}
  \item observability can be context-dependent;
  \item not all questions are jointly representable;
  \item coarse-graining has formal information loss;
  \item distinguishability obeys monotonicity under channels;
  \item recoverability characterizes when apparent loss is
    reversible.
\end{itemize}
MONAD remains a separate mathematical question.

\subsection{Endogenous and recursive change}

MONAD must eventually support more than one level of description: a
hierarchy $\Ctx_0, \Ctx_1, \Ctx_2, \ldots$ in which the rules defining
admissible transformations at one level become objects at the next.
Impredicative self-reference is not needed in the first theory. A
theory confined to one fixed description layer is incomplete relative
to the research goal. The recursive extension is a later theorem
target.

\subsubsection{Higher-order change}

MONAD distinguishes change of an entity ($x \to x'$), change
of the description ($C \to C'$), and change of the rules for changing
descriptions ($\TransportOp_C \to \TransportOp_{C'}$). The third
change modifies the transformation space itself.

\subsubsection{The self-similar pattern}

At many levels the same pattern recurs: description, compression or
transformation, residual defect, description refinement, new
transformation. This is structural self-similarity, not a claim of
literal fractal geometry. The theory must only permit the
transformation problem to recur on transformed description spaces.

\subsubsection{Confusion as a valid state}

When different members of a current equivalence class can produce
different outcomes under admissible but unresolved extensions, the
current description does not support a unique stronger statement. The
valid result is \emph{underdetermined}. This protects the theory from
false certainty.

\subsubsection{The representation-residual principle}

If a representation explains part of a structure and leaves a residual
$\varepsilon$, the theory must not automatically classify
$\varepsilon$ as noise. The residual can indicate measurement noise,
model error, an omitted variable, a missing distinction, a missing
dimension, or the wrong description space. Criteria separating these
possibilities are a major open problem. A theory that always absorbs
residuals into parameter error cannot satisfy generative openness.

\subsection{Proved baselines and remaining questions}

The exact set-based results are the observational equivalence,
coarsest quotient, unique descent, family-growth, universal compression,
empty-defect and covered-extension theorems above. Their precise scopes
are stated where proved; a result in this substrate does not settle
every enriched or resource-bounded version of the corresponding question.

The remaining research questions are:
\begin{description}
  \item[Resource minimality.] Which objective is minimized: class count,
    bits, computation, update, measurement, communication or decision loss?
  \item[Recovery.] What information remains accessible, which recovery
    operations are admitted, and what is their minimum cost?
  \item[Composition.] When does accounting compose exactly, and what
    demonstrated obstruction remains when it does not?
  \item[Path dependence.] When is endpoint-only transport sufficient?
    Do differences persist after ordinary hidden state is restored?
  \item[Trajectory relevance.] How does a changed policy or trajectory
    family change the distinctions a present representation must preserve?
  \item[Expressivity.] Did selection, budget, probes, language, operations,
    carrier or access change? Non-factorization alone does not answer this.
  \item[Substrate.] What minimal structure supports the requested operations?
  \item[Endogenous and recursive change.] What constructs, admits and
    validates a new context, and when can transformations become objects
    without uncontrolled self-reference?
\end{description}

\subsubsection{Approximate and resource-bounded MONAD}

The exact set-based theory is the starting point. The resource and
statistical specializations in the following sections study parts of the frontier
$\operatorname{Defect}(q_B)$ against a resource budget $B$, with
bounds of the form
$\inf_{\operatorname{cost}(q) \le B} \operatorname{Risk}(q)$. Cost and
risk depend on the chosen substrate. A general law relating these
quantities across changing contexts remains open.

\subsubsection{A possible probabilistic anticipation layer}

A probabilistic layer can place a measure over continuations and
derive relevance from path-space expectation, with the scalar form
$R_C(d) = \mathbb{E}_{\gamma \sim \mu}[U(d, \gamma)]$ only a
specialization. This layer is also not part of the minimal pure
theory.

\subsubsection{The determinant criterion}

A mature theory makes these facts intrinsic, refining the signature
of \Cref{ss:signature}:
\begin{itemize}
  \item whether a transformation collapses a distinction;
  \item whether it creates one;
  \item whether apparent loss is recoverable;
  \item whether an invariant survives;
  \item whether path matters;
  \item whether the present basis is sufficient;
  \item whether a future extension invalidates a prior compression.
\end{itemize}
The target is not a specific formula. The target is the correct
object.

\subsubsection{Correspondence requirements}

These are required correspondences, not a list of completed reductions.
For fixed probes, the set-based construction gives observational
equivalence; a behavioral or bisimulation correspondence additionally
needs the transition assumptions of \Cref{sss:coalgebra}. The exact
partition abstraction correspondence is proved in
\Cref{thm:partition-completeness,cor:budget-completeness}.
Probabilistic specializations must match the relevant data-processing
and sufficiency results; the quantum comparison in \Cref{ss:qit}
requires its own state and channel structure. A smooth-bundle
specialization with a connection should recover ordinary transport and
holonomy (\Cref{ss:holonomy}); an exact-category specialization should
use its kernel, image, cokernel and exact-sequence machinery
(\Cref{ss:homological}). The latter specializations remain targets.
MONAD-FORC~0.2
(\Cref{sec:state-space} onward) is the first fixed-context
correspondence target of this programme.

\subsubsection{Separation from applications}

MONAD is the mathematics of changing descriptions and
distinguishability. A governed operational system can freeze a
description and operate within it. Application and research systems
can use MONAD structure, test it, and supply examples. They must not
define its axioms.

\subsection{The complete MONAD problem}

Let $\Ctx$ be a compositional family of description contexts with
carriers, probe families, and continuation-induced
$\contEq{C}$. MONAD seeks a minimal, representation-independent
calculus that does all of the following:
\begin{enumerate}
  \item \textbf{It represents generative openness.} It identifies when
    a context extension introduces distinctions that cannot factor
    through the prior representation.
  \item \textbf{It defines anticipatory compression.} It identifies
    the coarsest representation that preserves the distinctions a
    continuation family requires, and characterizes how this fails
    when the family expands.
  \item \textbf{It gives invariant accounting.} It distinguishes
    preservation, temporary inaccessibility, recoverability,
    irreversible collapse, and novelty, with a compositional law or an
    explicit obstruction.
  \item \textbf{It gives trajectory-sensitive relevance.} Relevance need
    not be scalar; it depends on history and admissible continuations and
    detects required distinctions the representation cannot express.
  \item \textbf{It enforces compatibility.} The four structures
    constrain one another through theorems, and unknown distinctions
    stay unresolved until the theory justifies their treatment.
  \item \textbf{It exposes the limits of its own substrate.} Structure
    is added only when a theorem, counterexample, or impossibility
    result forces the addition.
  \item \textbf{It supports recursive extension.} Description
    transformations become objects at a higher level without letting a
    description validate its own change rules.
\end{enumerate}

\Principle{MONAD asks: how can a mathematical system define, compress,
transport, and revise its own space of distinctions when admissible
future continuations, historical trajectories, and transformations can
change what is distinguishable?}

The four primary requirements are generative openness, anticipatory
compression, invariant accounting, and trajectory-sensitive relevance.
The higher-order requirement is compatibility and coherence.
\begin{itemize}
  \item The tension: finite compression against future
    distinctions that are not yet representable.
  \item The epistemic rule: indistinguishable now is not
    identical forever.
  \item The obstruction: unrestricted openness with perfect
    future-safe compression forces trivial compression
    (\Cref{prop:no-go}).
  \item The objective: the intrinsic object that makes
    preservation, collapse, novelty, recoverability, and trajectory
    dependence directly visible.
\end{itemize}
